\section{Noether 正规化}\label{sec:noether-normalization}

前三节发展了若干模论概念. 现在, 我们将用这些概念研究处于代数几何核心的一类代数对象: 有限生成 $K$-代数.

Noether 正规化定理是关于所有有限生成 $K$-代数的结构定理. 为了说明这一定理所讨论的问题, 回想一下, 有限生成 $K$-代数的典型例子是多项式环
\[
  K[x_{1},\ldots,x_{d}].
\]
并非每个有限生成 $K$-代数都同构于多项式环. Noether 正规化定理试图回答: 用多项式环\textquotedbl{}逼近\textquotedbl{}有限生成 $K$-代数, 究竟能逼近到什么程度? 我们将看到, 给定有限生成 $K$-代数 $A$, 总可以找到子代数 $B\subseteq A$, 使得
\begin{enumerate}
  \item 对某个 $d$, 有 $B\cong K[x_{1},\ldots,x_{d}]$;
  \item $A$ 是有限生成 $B$-模, 或者等价地, $A$ 在 $B$ 上整.
\end{enumerate}
换句话说, Noether 正规化定理保证, 每个有限生成 $K$-代数都是某个多项式环 (最简单的一类 $K$-代数) 上的有限生成模 (最小的一类环扩张).

在陈述并证明 Noether 正规化定理之前, 我们先引入代数独立的概念. 它把线性独立推广到多项式代数的情形, 也便于我们讨论判定子代数何时同构于多项式环的准则.

\begin{Definition}{代数相关与代数独立}\leavevmode\label{def:algebraic-dependence-and-independence}\strut\par\nopagebreak[4]\noindent
设 $A$ 是 $R$-代数. 若存在非零多项式 $f\in R[x_{1},\ldots,x_{d}]$, 使得
\[
  f(a_{1},\ldots,a_{d})=0,
\]
则称元素 $a_{1},\ldots,a_{d}\in A$ 在 $R$ 上\emph{代数相关}. 如果不存在这样的多项式, 则称 $a_{1},\ldots,a_{d}$ 在 $R$ 上\emph{代数独立}.

若集合 $S\subseteq A$ 含有互不相同且在 $R$ 上代数相关的元素 $a_{1},\ldots,a_{d}$, 则称 $S$ 在 $R$ 上\emph{代数相关}. 如果不存在这样的元素, 则称 $S$ 在 $R$ 上\emph{代数独立}.
\end{Definition}

这里的用语有一个细微之处. 元素的列表允许重复, 而集合不计重复. 因此, 即便集合 $\{a_{1},\ldots,a_{d}\}$ 并不代数相关, 元素 $a_{1},\ldots,a_{d}$ 仍可能代数相关. 这种情况恰好发生在这些元素有重复的时候. 例如, 若 $a_{1}=a_{2}=\pi$, 元素 $a_{1}$ 和 $a_{2}$ 在 $\mathbb{Z}$ 上代数相关, 因为非零多项式 $f=x_{1}-x_{2}\in\mathbb{Z}[x_{1},x_{2}]$ 在 $(a_{1},a_{2})$ 处取值为零. 但集合 $\{a_{1},a_{2}\}=\{\pi\}$ 在 $\mathbb{Z}$ 上代数独立. 为避免混淆, 我们会在必要时强调 $a_{1},\ldots,a_{d}$ 互不相同.

由有限个代数独立元素生成的子代数同构于多项式环, 如下一个结论所述. 这是第一同构定理的一个应用 (Exercise~\ref{ex:prove-independent-generators-proposition}).

\begin{Proposition}{代数独立的生成元}\leavevmode\label{prop:algebraically-independent-generators}
设 $A$ 是 $R$-代数. 元素 $a_{1},\ldots,a_{d}\in A$ 在 $R$ 上代数独立, 当且仅当取值映射
\[
  \begin{aligned}
    R[x_{1},\ldots,x_{d}]&\to R[a_{1},\ldots,a_{d}]\subseteq A,\
    f(x_{1},\ldots,x_{d})&\mapsto f(a_{1},\ldots,a_{d})
  \end{aligned}
\]
是 $R$-代数同构.
\end{Proposition}

换句话说, Proposition~\ref{prop:algebraically-independent-generators} 断言, 给定 $R$-代数 $A$, 找到一个同构于多项式环的子代数 $B$, 等价于找到一个由代数独立元素组成的有限集合. 由此引出下面的 Noether 正规化定理, 它对任意域 $K$ 都成立.

\begin{Theorem}{Noether 正规化定理}\leavevmode\label{thm:noether-normalization}
若 $A$ 是有限生成 $K$-代数, 则存在有限子集 $S\subseteq A$, 使得
\begin{enumerate}
  \item[(i)] $S$ 在 $K$ 上代数独立;
  \item[(ii)] $A$ 在 $K[S]$ 上整.
\end{enumerate}
\end{Theorem}

具体地说, 若将定理中集合 $S$ 的各个不同元素写成 $a_{1},\ldots,a_{d}$, 则 $K[a_{1},\ldots,a_{d}]$ 同构于 $d$ 元多项式环, 且 $A$ 在 $K[a_{1},\ldots,a_{d}]$ 上整. 我们未必能找到一个代数独立的有限集合, 使它生成整个 $K$-代数 $A$ (因为 $A$ 未必同构于多项式环). 然而, $A$ 在 $K[a_{1},\ldots,a_{d}]$ 上整, 因而是有限生成 $K[a_{1},\ldots,a_{d}]$-模. 这表明 $a_{1},\ldots,a_{d}$ 所生成的多项式环\textquotedbl{}尽可能接近 $A$\textquotedbl{}.

Noether 正规化定理启发了下面的定义.

\begin{Definition}{Noether 基}\leavevmode\label{def:noether-basis}
设 $A$ 是有限生成 $K$-代数. $A$ 在 $K$ 上的一个\emph{Noether 基}是满足下列条件的有限子集 $S\subseteq A$:
\begin{enumerate}
  \item[(i)] $S$ 在 $K$ 上代数独立;
  \item[(ii)] $A$ 在 $K[S]$ 上整.
\end{enumerate}
\end{Definition}

用 Noether 基的语言来说, Noether 正规化定理就是断言每个有限生成 $K$-代数都存在 Noether 基. 在深入证明这一定理之前, 先来看几个具体的 Noether 基.

\begin{Example}{{\textmd{$K[x,y]/\langle x^{2}+y^{2}-1\rangle$} 的一个 Noether 基}}\leavevmode\label{ex:noether-basis-for-circle}
考虑 $K$-代数
\[
  A=\frac{K[x,y]}{\langle x^{2}+y^{2}-1\rangle}.
\]
注意, $A$ 由 $[x]$ 和 $[y]$ 生成. 但这两个元素并不代数独立, 因为它们满足一个非零的多项式关系:
\[
  [x]^{2}+[y]^{2}-1=[x^{2}+y^{2}-1]=0.
\]
令 $a=[x]\in A$. 我们证明 $\{a\}$ 是 $A$ 的一个 Noether 基.

首先, $a$ 在 $K$ 上代数独立. 为证明这一点, 必须说明对任意非零一元多项式 $f\in K[z]$, 有 $f(a)\neq0$. 注意任意这样的 $f$ 都满足 $f(x)\notin\langle x^{2}+y^{2}-1\rangle$. 因此,
\[
  f(a)=[f(x)]\neq0\in\frac{K[x,y]}{\langle x^{2}+y^{2}-1\rangle}.
\]
其次, 注意到 $K[a]\neq A$, 因为我们无法将 $[y]\in A$ 写成 $a=[x]$ 的多项式表达式. 然而, Example~\ref{ex:circle-coordinate-ring-module} 已经说明, 作为 $K[a]$-模, $A$ 由 $[1]$ 和 $[y]$ 生成:
\[
  A=K[a]\{[1],[y]\}.
\]
这就证明 $A$ 是有限生成 $K[a]$-模, 因此它在 $K[a]$ 上整. 由于 $a$ 在 $K$ 上代数独立, 而 $A$ 在 $K[a]$ 上整, 所以 $\{a\}$ 是一个 Noether 基.

这个例子容易推广到单位 $n$-球面的坐标环
\[
  A=\frac{K[x_{1},\ldots,x_{n}]}{\langle x_{1}^{2}+\cdots+x_{n}^{2}-1\rangle},
\]
从而说明 $\{a_{1}=[x_{1}],\ldots,a_{n-1}=[x_{n-1}]\}$ 是一个 Noether 基 (Exercise~\ref{ex:noether-basis-for-unit-sphere}).
\end{Example}

\begin{Example}{{\textmd{$K[x,y]/\langle xy-1\rangle$} 的一个 Noether 基}}\leavevmode\label{ex:noether-basis-for-hyperbola}
考虑 $K$-代数
\[
  A=\frac{K[x,y]}{\langle xy-1\rangle}.
\]
和 Example~\ref{ex:noether-basis-for-circle} 一样, $A$ 由 $[x]$ 和 $[y]$ 生成, 但这两个元素并不代数独立, 因而不构成 Noether 基. 自然会猜想, 单独取 $[x]$ 或 $[y]$ 也许就能构成 Noether 基. 但事实并非如此. 例如, $[x]$ 固然在 $K$ 上代数独立, 却可以验证 $A$ 不在 $K[x]$ 上整 (Exercise~\ref{ex:hyperbola-not-integral-over-x}).

尽管单独取 $[x]$ 或 $[y]$ 都不能构成 $A$ 的 Noether 基, 但 Noether 正规化定理保证 Noether 基一定存在. 在本例中, 一个 Noether 基由元素 $a=[x+y]$ 给出. 证明的思路见 Exercise~\ref{ex:hyperbola-integral-over-sum}.
\end{Example}

看过几个 Noether 基的例子后, 现在转向 Noether 正规化定理的证明. 证明采用一个巧妙的归纳论证, 其中关键是 Theorem~\ref{thm:finite-generation-integrality}.

\begin{Proof}\textbf{of Theorem~\ref{thm:noether-normalization}.}\quad
按定义, 任意有限生成 $K$-代数都可写成 $A=K[b_{1},\ldots,b_{n}]$, 其中 $b_{1},\ldots,b_{n}\in A$. 我们对生成元的个数 $n$ 归纳.

(归纳起点) 若 $n=0$, 则 $A=K$, 空集就是一个 Noether 基.

(归纳步骤) 假设生成元少于 $n$ 个的所有 $K$-代数都有 Noether 基, 并令 $A=K[b_{1},\ldots,b_{n}]$. 如果 $b_{1},\ldots,b_{n}$ 在 $K$ 上代数独立, 那么 $\{b_{1},\ldots,b_{n}\}$ 就是 $A$ 在 $K$ 上的一个 Noether 基, 结论得证. 于是设 $b_{1},\ldots,b_{n}$ 在 $K$ 上不代数独立, 即存在非零多项式 $f\in K[x_{1},\ldots,x_{n}]$, 使得
\[
  f(b_{1},\ldots,b_{n})=0.
\]
我们将变换这个关系, 得到一个关于 $b_{1}$ 首一的关系.

选取整数 $N$, 使之大于 $f$ 中所有非零系数单项式里出现的最大指数. 定义新元素 $c_{2},\ldots,c_{n}\in A$ 如下:
\[
  c_{2}=b_{2}-b_{1}^{N},\quad c_{3}=b_{3}-b_{1}^{N^{2}},\quad\ldots,\quad
  c_{n}=b_{n}-b_{1}^{N^{n-1}},
\]
因而
\[
  0=f(b_{1},\ldots,b_{n})
  =f\bigl(b_{1},c_{2}+b_{1}^{N},c_{3}+b_{1}^{N^{2}},\ldots,c_{n}+b_{1}^{N^{n-1}}\bigr).
\]
考虑下面的一元多项式
\[
  g(z)=f\bigl(z,c_{2}+z^{N},c_{3}+z^{N^{2}},\ldots,c_{n}+z^{N^{n-1}}\bigr)
  \in K[c_{2},\ldots,c_{n}][z].
\]
注意 $g(b_{1})=0$. 接下来证明 $g$ 的首项系数是单位, 从而 $b_{1}$ 在 $K[c_{2},\ldots,c_{n}]$ 上整.

将 $f$ 的任一单项式 $x_{1}^{\alpha_{1}}\cdots x_{n}^{\alpha_{n}}$ 在 $\bigl(z,c_{2}+z^{N},\ldots,c_{n}+z^{N^{n-1}}\bigr)$ 处取值, 得到
\begin{equation}\label{eq:noether-normalization-evaluated-monomial}
  z^{\alpha_{1}}(c_{2}+z^{N})^{\alpha_{2}}(c_{3}+z^{N^{2}})^{\alpha_{3}}
    \cdots(c_{n}+z^{N^{n-1}})^{\alpha_{n}}.
\end{equation}
合并 $z$ 的同次幂, 可将 \eqref{eq:noether-normalization-evaluated-monomial} 写成
\[
  z^{\alpha_{1}+\alpha_{2}N+\alpha_{3}N^{2}+\cdots+\alpha_{n}N^{n-1}}
  +\text{$z$ 的较低次项}.
\]
由于选取的 $N$ 大于 $\alpha_{1},\ldots,\alpha_{n}$, 指数
\[
  \alpha_{1}+\alpha_{2}N+\alpha_{3}N^{2}+\cdots+\alpha_{n}N^{n-1}
\]
唯一确定 $\alpha_{1},\ldots,\alpha_{n}$. (这个断言就是\textquotedbl{}$N$ 进制展开的唯一性\textquotedbl{}: 例如, 若 $N=10$, 它就是断言一个数的各位数字由这个数本身唯一确定. 参见 Exercise~\ref{ex:uniqueness-of-base-n-expansions}.) 特别地, 对不同的单项式
\[
  x_{1}^{\alpha_{1}}\cdots x_{n}^{\alpha_{n}}\qquad\text{和}\qquad
  x_{1}^{\alpha'_{1}}\cdots x_{n}^{\alpha'_{n}},
\]
在 $\bigl(z,c_{2}+z^{N},\ldots,c_{n}+z^{N^{n-1}}\bigr)$ 处取值后, 得到的 $z$ 的次数不同.

因此, $f$ 中有唯一一个非零系数的单项式使 $\alpha_{1}+\alpha_{2}N+\alpha_{3}N^{2}+\cdots+\alpha_{n}N^{n-1}$ 达到最大. 若 $c\in K$ 是该单项式在 $f$ 中的非零系数, 则 $c$ 也是 $g$ 的首项系数. 所以 $c^{-1}g\in K[c_{2},\ldots,c_{n}][z]$ 是一个在 $b_{1}$ 处取值为零的首一多项式, 因而 $b_{1}$ 在 $K[c_{2},\ldots,c_{n}]$ 上整. 由 Theorem~\ref{thm:finite-generation-integrality} 得到 $K[c_{2},\ldots,c_{n}][b_{1}]$ 在 $K[c_{2},\ldots,c_{n}]$ 上整. 又由于 $b_{1},b_{2},\ldots,b_{n}$ 的多项式组合和 $b_{1},c_{2},\ldots,c_{n}$ 的多项式组合可以相互转换, 有
\[
  K[c_{2},\ldots,c_{n}][b_{1}]=K[b_{1},c_{2},\ldots,c_{n}]
    =K[b_{1},b_{2},\ldots,b_{n}]=A.
\]
于是证明了 $A=K[b_{1},\ldots,b_{n}]$ 在 $K[c_{2},\ldots,c_{n}]$ 上整.

由归纳假设, 存在一个 Noether 基 $S\subseteq K[c_{2},\ldots,c_{n}]$. 我们断言 $S$ 也是 $A$ 的一个 Noether 基. 为此, 需要证明 $S$ 在 $K$ 上代数独立, 且 $A$ 在 $K[S]$ 上整. $S$ 在 $K$ 上代数独立直接来自 $S$ 是 $K[c_{1},\ldots,c_{n}]$ 的一个 Noether 基这一事实. 为证明 $A$ 在 $K[S]$ 上整, 注意下列各环都在它前一个环上整:
\[
  K[S]\subseteq K[c_{2},\ldots,c_{n}]\subseteq A.
\]
其中, 第一个包含关系的整性来自 $S$ 是 $K[c_{1},\ldots,c_{n}]$ 的一个 Noether 基这一假设, 第二个包含关系的整性则是前面论证的核心. 由 Theorem~\ref{thm:finite-generation-integrality}, 这些环中的每一个作为前一个环上的模都是有限生成的. 因此由 Exercise~\ref{ex:finite-module-in-tower} 得到 $A$ 作为 $K[S]$-模有限生成; 再用 Theorem~\ref{thm:finite-generation-integrality}, 得到 $A$ 在 $K[S]$ 上整. 所以 $S$ 是 $A$ 的一个 Noether 基, 归纳论证完成.
\end{Proof}

Noether 正规化定理是过去四节代数理论发展的结晶, 这些努力将带来许多收获. 最直接的一项是, 如下一节所见, Noether 正规化可用来证明零点定理, 从而证明代数与几何的等价性. 收获还不止于此: Noether 正规化也与维数密切相关, 这是下一章要引入的概念. 我们将看到, 若 $X$ 是仿射簇, 则坐标环 $K[X]$ 的任意 Noether 基所含元素的个数等于 $X$ 的维数. 借此可以用 Noether 基证明维数的基本性质.

\BookExerciseHeading
\begin{BookExercises}
  \item\label{ex:prove-independent-generators-proposition} 证明 Proposition~\ref{prop:algebraically-independent-generators}.
  \item\label{ex:noether-basis-for-unit-sphere} 考虑 $K$-代数
  \[
    A=\frac{K[x_{1},\ldots,x_{n}]}{\langle x_{1}^{2}+\cdots+x_{n}^{2}-1\rangle}.
  \]
  对所有 $i=1,\ldots,n-1$, 令 $a_{i}=[x_{i}]$.
  \begin{enumerate}
    \item 证明 $a_{1},\ldots,a_{n-1}$ 代数独立.
    \item 证明 $A$ 在 $K[a_{1},\ldots,a_{n-1}]$ 上整.
  \end{enumerate}
  \item\label{ex:hyperbola-not-integral-over-x} 令 $A=K[x,y]/\langle xy-1\rangle$, 并令 $a=[x]$.
  \begin{enumerate}
    \item 证明 $a$ 在 $K$ 上代数独立.
    \item 证明 $A$ 不在 $K[a]$ 上整.
  \end{enumerate}
  \item\label{ex:hyperbola-integral-over-sum} 令 $A=K[x,y]/\langle xy-1\rangle$, 并令 $a=[x+y]$.
  \begin{enumerate}
    \item 证明 $a$ 在 $K$ 上代数独立.
    \item 证明 $A$ 在 $K[a]$ 上整. (提示: 证明 $[x]$ 和 $[y]$ 都是 $K[a]$ 上首一二次多项式的零点.)
  \end{enumerate}
  \item\label{ex:noether-basis-is-maximal-independent} 设 $A$ 是 $K$-代数. 证明 $A$ 在 $K$ 上的每个 Noether 基都是极大的代数独立集. (换句话说, 真包含一个 Noether 基的任何集合都必定代数相关.)
  \item\label{ex:maximal-independent-need-not-be-noether-basis} 举例说明 Exercise~\ref{ex:noether-basis-is-maximal-independent} 的逆命题不成立. 也就是说, 若 $A$ 是有限生成 $K$-代数, 且 $S$ 是极大的代数独立集, $A$ 不一定在 $K[S]$ 上整.
  \item\label{ex:uniqueness-of-base-n-expansions} 设 $N$ 是正整数, 令 $[N]=\{0,1,\ldots,N-1\}$. 对 $k\geq0$, 考虑函数
  \[
    \begin{aligned}
      \psi_{n}:[N]^{k+1} &\to \mathbb{N},\\
      (\alpha_{0},\ldots,\alpha_{k}) &\mapsto
        \alpha_{0}+\alpha_{1}N+\alpha_{2}N^{2}+\cdots+\alpha_{k}N^{k}.
    \end{aligned}
  \]
  \begin{enumerate}
    \item 若 $N=10$, 你会怎样描述数 $\psi_{k}(\alpha_{0},\ldots,\alpha_{k})$?
    \item 证明对所有 $k\geq0$, $\psi_{k}$ 都是单射. (换句话说, 数 $\alpha_{0}+\alpha_{1}N+\alpha_{2}N^{2}+\cdots+\alpha_{k}N^{k}$ 唯一确定 $\alpha_{0},\ldots,\alpha_{k}$.)
  \end{enumerate}
\end{BookExercises}
